\documentclass[11pt,a4paper]{ctexart}
\usepackage[margin=2.2cm]{geometry}
\usepackage{xcolor}
\usepackage{enumitem}
\usepackage{amssymb}
\usepackage[normalem]{ulem}
\usepackage{titlesec}
\usepackage{hyperref}
\definecolor{navy}{HTML}{1A365D}
\definecolor{ink}{HTML}{1F2933}
\definecolor{muted}{HTML}{66788A}
\definecolor{soft}{HTML}{EEF4FA}
\definecolor{gold}{HTML}{8A5A12}
\hypersetup{colorlinks=true,linkcolor=navy,urlcolor=navy}
\setlist{nosep}
\titleformat{\section}{\Large\bfseries\color{navy}}{}{0pt}{}
\titleformat{\subsection}{\large\bfseries\color{navy}}{}{0pt}{}
\pagestyle{plain}

\title{\textbf{七上冠词与代词专项}\\[4pt]
\large 人教版（新课纲版）七年级上册 · 冠词 / 代词}
\author{初中英语资料站}
\date{2026-11-12}
\begin{document}
\maketitle

\section*{使用说明}
\begin{itemize}[leftmargin=1.4em]
\item 本讲义对应人教版（新课纲版）七年级上册语法要求：冠词与代词。
\item 共 4 讲语法、6 级 48 题。
\item 在线版：\url{https://1901151458-prog.github.io/english-materials/g7a-articles/}
\end{itemize}

\section{定冠词 the}
定冠词的基本概念是「\textbf{特指}」，用于指说话人和听话人\textbf{已知}的人或事物，可与单数可数名词、复数可数名词或不可数名词一起使用。
\begin{itemize}[leftmargin=1.4em]
\item Where is \underline{the} library?
\item \underline{The} glasses are on her dad's nose.
\item \underline{The} milk is on the table.
\end{itemize}
\textbf{读音：}辅音前读 /ðə/，元音前读 /ði/，特别强调时可读 /ðiː/。

\section{不定冠词 a / an}
不定冠词的基本概念是「\textbf{非特指}」，指某类人或物中的任何一个或某一个，\textbf{只能与单数可数名词}一起使用。
\begin{itemize}[leftmargin=1.4em]
\item I have \underline{a} pet dog.（该类中的一例）
\item Sometimes I play basketball for \underline{an} hour.（一个小时）
\item We have three lessons \underline{a} day.（每一个）
\end{itemize}
\textbf{an} 放在以元音（音素）开头的名词前面，如 an eraser、an hour（h 不发音）。

\section{零冠词}
零冠词指名词前\textbf{不加冠词}的情况。
\begin{itemize}[leftmargin=1.4em]
\item My favourite subject is \underline{English}.
\item Binbin is good with \underline{numbers}.
\item It's time for me to go to \underline{bed}.
\end{itemize}
\textbf{教材注释：}有的语法书并不把零冠词列为冠词中的一种。

\section{代词}
\begin{center}
\begin{tabular}{p{2.6cm}p{3.0cm}p{3.0cm}p{4.0cm}}
\hline
\textbf{人称} & \textbf{主格} & \textbf{宾格} & \textbf{形容词性物主}\\
\hline
一单 & I & me & my\\
二单 & you & you & your\\
三单 & he / she / it & him / her / it & his / her / its\\
一复 & we & us & our\\
三复 & they & them & their\\
\hline
\end{tabular}
\end{center}
\textbf{指示代词：}this / these 指时间或空间上\textbf{较近}的事物；that / those 指\textbf{较远}的。


\newpage
\section{分级练习（48 题）}
\noindent 按难度分级排列，先做基础，再做成分与扩展，最后做综合题。

\subsection*{第 1 级 · 定冠词 the}
\noindent\textcolor{gray}{特指双方已知的人或事物；/ðə/ 与 /ði/。}\par\smallskip
\begin{enumerate}[leftmargin=*,itemsep=4pt]
\item Where is \_\_\_ library?
\begin{enumerate}[label=\Alph*.,leftmargin=2em,itemsep=1pt,topsep=2pt]
\item the
\item a
\item an
\item 不填
\end{enumerate}
\item \_\_\_ glasses are on her dad’s nose.
\begin{enumerate}[label=\Alph*.,leftmargin=2em,itemsep=1pt,topsep=2pt]
\item The
\item A
\item An
\item 不填
\end{enumerate}
\item \_\_\_ milk is on the table.
\begin{enumerate}[label=\Alph*.,leftmargin=2em,itemsep=1pt,topsep=2pt]
\item The
\item A
\item An
\item 不填
\end{enumerate}
\item 定冠词的基本概念是 \_\_\_。
\begin{enumerate}[label=\Alph*.,leftmargin=2em,itemsep=1pt,topsep=2pt]
\item 特指
\item 非特指
\item 泛指
\item 任指
\end{enumerate}
\item 定冠词在辅音前读 \_\_\_。
\begin{enumerate}[label=\Alph*.,leftmargin=2em,itemsep=1pt,topsep=2pt]
\item /ðə/
\item /ði/
\item /θə/
\item /θi/
\end{enumerate}
\item 定冠词在元音前读 \_\_\_。
\begin{enumerate}[label=\Alph*.,leftmargin=2em,itemsep=1pt,topsep=2pt]
\item /ðə/
\item /ði/
\item /θə/
\item /zi/
\end{enumerate}
\item \_\_\_ sun rises in the east.
\begin{enumerate}[label=\Alph*.,leftmargin=2em,itemsep=1pt,topsep=2pt]
\item The
\item A
\item An
\item 不填
\end{enumerate}
\item 定冠词可与 \_\_\_ 一起使用。
\begin{enumerate}[label=\Alph*.,leftmargin=2em,itemsep=1pt,topsep=2pt]
\item 单数可数、复数可数和不可数名词
\item 只能单数可数名词
\item 只能复数名词
\item 只能不可数名词
\end{enumerate}
\end{enumerate}

\subsection*{第 2 级 · 不定冠词 a / an}
\noindent\textcolor{gray}{非特指；an 用于元音音素开头。}\par\smallskip
\begin{enumerate}[leftmargin=*,itemsep=4pt]
\item I have \_\_\_ pet dog.
\begin{enumerate}[label=\Alph*.,leftmargin=2em,itemsep=1pt,topsep=2pt]
\item a
\item an
\item the
\item 不填
\end{enumerate}
\item Sometimes I play basketball for \_\_\_ hour.
\begin{enumerate}[label=\Alph*.,leftmargin=2em,itemsep=1pt,topsep=2pt]
\item a
\item an
\item the
\item 不填
\end{enumerate}
\item We have three lessons \_\_\_ day.
\begin{enumerate}[label=\Alph*.,leftmargin=2em,itemsep=1pt,topsep=2pt]
\item a
\item an
\item the
\item 不填
\end{enumerate}
\item an 放在以 \_\_\_ 开头的名词前面。
\begin{enumerate}[label=\Alph*.,leftmargin=2em,itemsep=1pt,topsep=2pt]
\item 元音音素
\item 辅音音素
\item 元音字母
\item 辅音字母
\end{enumerate}
\item \_\_\_ eraser
\begin{enumerate}[label=\Alph*.,leftmargin=2em,itemsep=1pt,topsep=2pt]
\item a
\item an
\item the
\item 不填
\end{enumerate}
\item \_\_\_ book
\begin{enumerate}[label=\Alph*.,leftmargin=2em,itemsep=1pt,topsep=2pt]
\item a
\item an
\item the
\item 不填
\end{enumerate}
\item 不定冠词只能与 \_\_\_ 一起使用。
\begin{enumerate}[label=\Alph*.,leftmargin=2em,itemsep=1pt,topsep=2pt]
\item 单数可数名词
\item 复数可数名词
\item 不可数名词
\item 任何名词
\end{enumerate}
\item 不定冠词的基本含义有「该类中的一例」「只有一个」「\_\_\_」。
\begin{enumerate}[label=\Alph*.,leftmargin=2em,itemsep=1pt,topsep=2pt]
\item 每一个
\item 全部
\item 某些
\item 任何一个
\end{enumerate}
\end{enumerate}

\subsection*{第 3 级 · 零冠词}
\noindent\textcolor{gray}{名词前不加冠词的情况。}\par\smallskip
\begin{enumerate}[leftmargin=*,itemsep=4pt]
\item My favourite subject is \_\_\_.
\begin{enumerate}[label=\Alph*.,leftmargin=2em,itemsep=1pt,topsep=2pt]
\item English
\item the English
\item an English
\item a English
\end{enumerate}
\item It’s time for me to go to \_\_\_.
\begin{enumerate}[label=\Alph*.,leftmargin=2em,itemsep=1pt,topsep=2pt]
\item bed
\item the bed
\item a bed
\item an bed
\end{enumerate}
\item Binbin is good with \_\_\_.
\begin{enumerate}[label=\Alph*.,leftmargin=2em,itemsep=1pt,topsep=2pt]
\item numbers
\item the numbers
\item a numbers
\item an numbers
\end{enumerate}
\item 零冠词指名词前 \_\_\_ 的情况。
\begin{enumerate}[label=\Alph*.,leftmargin=2em,itemsep=1pt,topsep=2pt]
\item 不加冠词
\item 加 the
\item 加 a
\item 加 an
\end{enumerate}
\item 有的语法书 \_\_\_ 把零冠词列为冠词的一种。
\begin{enumerate}[label=\Alph*.,leftmargin=2em,itemsep=1pt,topsep=2pt]
\item 并不
\item 也会
\item 必须
\item 一定
\end{enumerate}
\item 早饭前：\_\_\_。
\begin{enumerate}[label=\Alph*.,leftmargin=2em,itemsep=1pt,topsep=2pt]
\item before breakfast
\item before the breakfast
\item before a breakfast
\item before an breakfast
\end{enumerate}
\item 球类运动前一般 \_\_\_。
\begin{enumerate}[label=\Alph*.,leftmargin=2em,itemsep=1pt,topsep=2pt]
\item 不加冠词
\item 加 the
\item 加 a
\item 加 an
\end{enumerate}
\item 在学校：\_\_\_。
\begin{enumerate}[label=\Alph*.,leftmargin=2em,itemsep=1pt,topsep=2pt]
\item at school
\item at the school
\item at a school
\item at an school
\end{enumerate}
\end{enumerate}

\subsection*{第 4 级 · 人称代词}
\noindent\textcolor{gray}{主格与宾格。}\par\smallskip
\begin{enumerate}[leftmargin=*,itemsep=4pt]
\item I 的宾格是 \_\_\_。
\begin{enumerate}[label=\Alph*.,leftmargin=2em,itemsep=1pt,topsep=2pt]
\item me
\item my
\item mine
\item I
\end{enumerate}
\item he 的宾格是 \_\_\_。
\begin{enumerate}[label=\Alph*.,leftmargin=2em,itemsep=1pt,topsep=2pt]
\item him
\item his
\item he
\item her
\end{enumerate}
\item they 的宾格是 \_\_\_。
\begin{enumerate}[label=\Alph*.,leftmargin=2em,itemsep=1pt,topsep=2pt]
\item them
\item their
\item they
\item theirs
\end{enumerate}
\item we 的宾格是 \_\_\_。
\begin{enumerate}[label=\Alph*.,leftmargin=2em,itemsep=1pt,topsep=2pt]
\item us
\item our
\item we
\item ours
\end{enumerate}
\item she 的宾格是 \_\_\_。
\begin{enumerate}[label=\Alph*.,leftmargin=2em,itemsep=1pt,topsep=2pt]
\item her
\item hers
\item she
\item he
\end{enumerate}
\item you 的宾格是 \_\_\_。
\begin{enumerate}[label=\Alph*.,leftmargin=2em,itemsep=1pt,topsep=2pt]
\item you
\item your
\item yours
\item you’s
\end{enumerate}
\item it 的宾格是 \_\_\_。
\begin{enumerate}[label=\Alph*.,leftmargin=2em,itemsep=1pt,topsep=2pt]
\item it
\item its
\item it’s
\item itself
\end{enumerate}
\item 「他们」的主格是 \_\_\_。
\begin{enumerate}[label=\Alph*.,leftmargin=2em,itemsep=1pt,topsep=2pt]
\item they
\item them
\item their
\item theirs
\end{enumerate}
\end{enumerate}

\subsection*{第 5 级 · 物主代词与指示代词}
\noindent\textcolor{gray}{形容词性物主代词；this/these/that/those。}\par\smallskip
\begin{enumerate}[leftmargin=*,itemsep=4pt]
\item I 的形容词性物主代词是 \_\_\_。
\begin{enumerate}[label=\Alph*.,leftmargin=2em,itemsep=1pt,topsep=2pt]
\item my
\item me
\item mine
\item I
\end{enumerate}
\item they 的形容词性物主代词是 \_\_\_。
\begin{enumerate}[label=\Alph*.,leftmargin=2em,itemsep=1pt,topsep=2pt]
\item their
\item them
\item they
\item theirs
\end{enumerate}
\item she 的形容词性物主代词是 \_\_\_。
\begin{enumerate}[label=\Alph*.,leftmargin=2em,itemsep=1pt,topsep=2pt]
\item her
\item hers
\item she
\item him
\end{enumerate}
\item \_\_\_ is my uncle’s farm.（这个）
\begin{enumerate}[label=\Alph*.,leftmargin=2em,itemsep=1pt,topsep=2pt]
\item This
\item That
\item These
\item Those
\end{enumerate}
\item \_\_\_ are my friends.（这些）
\begin{enumerate}[label=\Alph*.,leftmargin=2em,itemsep=1pt,topsep=2pt]
\item This
\item That
\item These
\item Those
\end{enumerate}
\item Do you see \_\_\_ red house by the river?（那个）
\begin{enumerate}[label=\Alph*.,leftmargin=2em,itemsep=1pt,topsep=2pt]
\item this
\item that
\item these
\item those
\end{enumerate}
\item \_\_\_ ping-pong bats in the corner are my grandpa’s.（那些）
\begin{enumerate}[label=\Alph*.,leftmargin=2em,itemsep=1pt,topsep=2pt]
\item This
\item That
\item These
\item Those
\end{enumerate}
\item this / these 指时间或空间上 \_\_\_ 的事物。
\begin{enumerate}[label=\Alph*.,leftmargin=2em,itemsep=1pt,topsep=2pt]
\item 较近
\item 较远
\item 中间
\item 两边
\end{enumerate}
\end{enumerate}

\subsection*{第 6 级 · 综合提升}
\noindent\textcolor{gray}{冠词、人称代词、物主代词与指示代词混合辨析。}\par\smallskip
\begin{enumerate}[leftmargin=*,itemsep=4pt]
\item \_\_\_ sun is shining brightly.
\begin{enumerate}[label=\Alph*.,leftmargin=2em,itemsep=1pt,topsep=2pt]
\item The
\item A
\item An
\item 不填
\end{enumerate}
\item I want to be \_\_\_ doctor.
\begin{enumerate}[label=\Alph*.,leftmargin=2em,itemsep=1pt,topsep=2pt]
\item a
\item an
\item the
\item 不填
\end{enumerate}
\item Look at \_\_\_ elephant!
\begin{enumerate}[label=\Alph*.,leftmargin=2em,itemsep=1pt,topsep=2pt]
\item a
\item an
\item the
\item 不填
\end{enumerate}
\item \_\_\_ is my schoolbag.（这个）
\begin{enumerate}[label=\Alph*.,leftmargin=2em,itemsep=1pt,topsep=2pt]
\item This
\item That
\item These
\item Those
\end{enumerate}
\item He is good at \_\_\_ English.
\begin{enumerate}[label=\Alph*.,leftmargin=2em,itemsep=1pt,topsep=2pt]
\item 不填
\item the
\item a
\item an
\end{enumerate}
\item \_\_\_ girl in red is my sister.
\begin{enumerate}[label=\Alph*.,leftmargin=2em,itemsep=1pt,topsep=2pt]
\item The
\item A
\item An
\item 不填
\end{enumerate}
\item I have \_\_\_ uncle and \_\_\_ aunt.
\begin{enumerate}[label=\Alph*.,leftmargin=2em,itemsep=1pt,topsep=2pt]
\item an; an
\item a; a
\item an; a
\item a; an
\end{enumerate}
\item that / those 指时间或空间上 \_\_\_ 的事物。
\begin{enumerate}[label=\Alph*.,leftmargin=2em,itemsep=1pt,topsep=2pt]
\item 较远
\item 较近
\item 中间
\item 两边
\end{enumerate}
\end{enumerate}

\newpage
\section*{答案与解析}
\small
\begin{enumerate}[leftmargin=1.6em,itemsep=2pt]
\item \textbf{A1-1 A}：双方都知道的那个图书馆，特指用 the；教材例句。
\item \textbf{A1-2 A}：特指那副眼镜，用 the；教材例句。
\item \textbf{A1-3 A}：特指桌上的牛奶，用 the；教材例句。
\item \textbf{A1-4 A}：教材：定冠词的基本概念是「特指」。
\item \textbf{A1-5 A}：教材：定冠词在辅音前读 /ðə/。
\item \textbf{A1-6 B}：教材：在元音前读 /ði/。
\item \textbf{A1-7 A}：世界上独一无二的事物前用 the。
\item \textbf{A1-8 A}：教材：可与单数可数名词、复数可数名词或不可数名词一起使用。
\item \textbf{A2-1 A}：pet 以辅音音素开头，用 a；教材例句。
\item \textbf{A2-2 B}：hour 以元音音素开头（h 不发音），用 an；教材例句。
\item \textbf{A2-3 A}：a day 表示「每一天」；教材例句。
\item \textbf{A2-4 A}：教材：an 放在以元音（音素）开头的名词前面。
\item \textbf{A2-5 B}：eraser 以元音音素开头，用 an；教材例句。
\item \textbf{A2-6 A}：book 以辅音音素开头，用 a。
\item \textbf{A2-7 A}：教材：不定冠词只能与单数可数名词一起使用。
\item \textbf{A2-8 A}：教材：其基本含义有「该类中的一例」「只有一个」「每一个」等。
\item \textbf{A3-1 A}：学科名前用零冠词；教材例句。
\item \textbf{A3-2 A}：go to bed 固定搭配，用零冠词；教材例句。
\item \textbf{A3-3 A}：泛指数字，用零冠词；教材例句。
\item \textbf{A3-4 A}：教材定义。
\item \textbf{A3-5 A}：教材注释：有的语法书并不把零冠词列为冠词中的一种。
\item \textbf{A3-6 A}：三餐名前用零冠词。
\item \textbf{A3-7 A}：play basketball 等球类运动前用零冠词。
\item \textbf{A3-8 A}：at school 是固定搭配，用零冠词。
\item \textbf{A4-1 A}：人称代词宾格；教材表3。
\item \textbf{A4-2 A}：教材表3。
\item \textbf{A4-3 A}：教材表3。
\item \textbf{A4-4 A}：教材表3。
\item \textbf{A4-5 A}：教材表3。
\item \textbf{A4-6 A}：you 的主格和宾格同形；教材表3。
\item \textbf{A4-7 A}：it 的主格和宾格同形；教材表3。
\item \textbf{A4-8 A}：教材表3。
\item \textbf{A5-1 A}：形容词性物主代词相当于形容词，置于名词前作定语。
\item \textbf{A5-2 A}：their children 等。
\item \textbf{A5-3 A}：her favourite subject 等；教材例句。
\item \textbf{A5-4 A}：this 指较近的事物；教材例句。
\item \textbf{A5-5 C}：these 是 this 的复数；教材例句。
\item \textbf{A5-6 B}：that 指较远的事物；教材例句。
\item \textbf{A5-7 D}：those 是 that 的复数；教材例句。
\item \textbf{A5-8 A}：教材表5。
\item \textbf{A6-1 A}：独一无二的事物前用 the。
\item \textbf{A6-2 A}：非特指的单数可数名词用 a。
\item \textbf{A6-3 B}：elephant 以元音音素开头，用 an。
\item \textbf{A6-4 A}：this 指近处单数事物。
\item \textbf{A6-5 A}：学科名前用零冠词。
\item \textbf{A6-6 A}：in red 限定了范围，特指用 the。
\item \textbf{A6-7 A}：uncle 和 aunt 都以元音音素开头，都用 an。
\item \textbf{A6-8 A}：教材表5。
\end{enumerate}

\end{document}
